\documentclass[11pt]{article}
\usepackage{fontspec}
\usepackage{amsmath,amssymb,amsthm,mathtools}
\usepackage[margin=1.1in]{geometry}
\usepackage{microtype}
\usepackage{parskip}
\usepackage{lmodern}
\usepackage[colorlinks=true,linkcolor=blue!50!black,urlcolor=blue!50!black]{hyperref}
\providecommand{\tightlist}{}
\title{Sealed-seam suppression in Froggatt–Nielsen-type models: two theorems, an obstruction by conjunction, and a discrete neutrino test}
\author{Felix Robles Elvira (ORCID: 0009-0009-2017-4394; independent researcher)}
\date{\small preprint, not peer reviewed, version 2026-06-11 (substantially rebuilt after hostile review; the withdrawn claims of the earlier draft are listed in Section 7 and stay withdrawn).}
\begin{document}
\maketitle

\section*{Abstract}
Froggatt–Nielsen (FN) mechanisms generate fermion mass hierarchies as powers of a small parameter $\lambda$, one insertion per unit of flavor charge, with $\lambda$ fitted.  We study an FN-type setting in which the suppression is \emph{pinned by a theorem}: for sealed two-outcome records — rank-one POVM-type letter pairs with weights $p$ and $1-p$ in a stationary representation — (i) the transition norm is $\sqrt{p(1-p)}\,|\langle u,v\rangle| \le \sqrt{p(1-p)}$, and (ii) a misaligned sealed pair cannot support a stationary state (top transfer eigenvalue $\tfrac12[1+\sqrt{1-4p(1-p) (1-|c|^2)}] < 1$ strictly), so stationarity forces saturation: $\lVert S\rVert = \sqrt{p(1-p)}$ exactly, with the seam phase pure rephasing gauge.  An exact selection rule follows in any algebra graded by a family charge whose only charged sealed generator is the seam: amplitudes vanish unless the seam count balances the charge difference (receipts: $0/27{,}000$ violations; bounds saturated to ten digits).  The paper's central \emph{negative} result is an \textbf{obstruction by conjunction}, found when the hostile review of an earlier draft demolished its seesaw realization and we completed the demolition: for \emph{integer} right-handed charges, Majorana diagonals carry even powers of $\lambda$ and an exhaustive Dirac-consistent scan over rule-allowed textures (charges in $\{0..3\}^3$, unit $\pm$ coefficients) never reaches the half-integer-rung spectrum (zero exact hits; best deviation $0.014$; receipted); for \emph{half-odd} right-handed charges the Majorana side can reach the rungs exactly, but the corresponding Dirac column carries half-integer charge, balances no integer seam count, and vanishes identically — rank collapse; and the textbook cancellation in any case removes RH charges from the light masses when the Dirac sector is charged consistently.  The surviving realization places half-unit charges on the lepton doublets (Weinberg-operator dressing), where the same selection rule yields a \textbf{discrete menu} for the light spectrum at unit coefficients (four full-rank members; two already excluded by data): the rung point $S = \sqrt{\varepsilon_{\rm eff}/(1+\varepsilon_{\rm eff})} = 0.172747$ (three insertions) and the cheaper texture $S = 0.167880$ (two insertions), with $\varepsilon_{\rm eff} = \varepsilon(1-\varepsilon)$ and $\varepsilon = 0.031768636446582$ a constant whose self-contained pedigree is the mathematical companion [C0].  Against JUNO's first data ($S = 0.17283 \pm 0.00167$, our combination of JUNO's $\Delta m^2_{21}$ with the global-fit $\Delta m^2_{31}$) the data-live members sit at $-0.05\sigma$ and $-3.0\sigma$; they are separated by $7.3\sigma$ at JUNO-final precision, so the menu is adjudicable this decade, and the minimality principle the program uses elsewhere favors — on current data — the \emph{worse} member, which we state rather than hide.  What this paper does \textbf{not} deliver, withdrawn from an earlier draft and listed with reasons: mixing-angle and CP predictions (they rested on an unlisted measure-zero orthogonality assumption), the claim that the mechanism's flavor symmetry uniquely "predicts its own mediator," and any charged-fermion statement.  One quantitative input is irreducible and stated in one line: the seam weight is the companion theory's first-order marginality, $p = \varepsilon$ — a referent scan shows it underivable from the obvious fixed-point quantities, and it carries the aggravating fact that the first-order law it comes from has its sign overturned at second order at the defining weight.  The look-elsewhere cost of the dressed value is counted (six of twelve natural one-factor forms fit today's central value; three are mutually inseparable even at JUNO-final), and the claim chronology relative to JUNO's first release is disclosed in full.

\section*{1. The Seam Theorem}
Work in the representation every reflection-positive stationary process admits (companion [C2], Thm 2.1): Hilbert space, PSD letters, $T = \sum_x A_x \preceq \mathbf 1$, stationary unit vector $T\omega = \omega$.

\textbf{Definition (S-atom premise, named).}  A \emph{sealed two-outcome seam} is a letter pair of rank one: $A_0 = p\,uu^\dagger$, $A_1 = (1-p)\,vv^\dagger$, unit $u, v$.

\textbf{Lemma 1 (norm).}  $\sqrt{A_0}\sqrt{A_1} = \sqrt{p(1-p)}\, \langle u,v\rangle\,uv^\dagger$, so $\lVert\sqrt{A_0}\sqrt{A_1} \rVert = \sqrt{p(1-p)}\,|\langle u,v\rangle| \le \sqrt{p(1-p)}$, with equality iff the poles align.

\textbf{Lemma 2 (stationarity gate).}  On the seam span, with $c = \langle u,v\rangle$ and $d = 4p(1-p)(1-|c|^2)$: $\lambda_{\max}(A_0+A_1) = \tfrac12[1 + \sqrt{1-d}\,]$, equal to $1$ iff $|c| = 1$; the exact gap is $\tfrac12[1 - \sqrt{1-d}\,] \ge p(1-p)(1-|c|^2)/2$, strictly positive for misaligned poles. (An earlier draft's abstract quoted the lower bound as the exact gap; corrected.)

\textbf{Theorem 1 (stationarity forces saturation).}  If the seam block acts invariantly at stationarity (premise: decoupling) and the stationary state overlaps it, then alignment is forced and the transition norm is exactly $\sqrt{p(1-p)}$.  At alignment the letters commute: the extremal seam is classical, and the factor is the two-pole weight of a definite record.  The seam phase ($u \mapsto e^{i\varphi}u$) changes no observable — the rephasing structure mass terms have always had.

\emph{Receipts}: norm identity to $2.8\times 10^{-16}$ (200 random complex seams, $d = 2..5$); eigenvalue closed form to $3.3\times 10^{-16}$.

\section*{2. The selection rule}
\textbf{Theorem 2.}  In an algebra carrying a graded charge $Q$ with every sealed generator $Q$-neutral except one seam $S$ shifting $Q$ by one unit, $\langle x'|W|x\rangle = 0$ for every word whose seam count fails to balance $x' - x$ (grade decomposition), and — since neutral letters are contractions — balanced amplitudes obey $|\langle x'|W|x\rangle| \le \lVert S\rVert^{|x'-x|}$ (submultiplicativity; an inequality, with saturation a property of the specific model, exhibited in the receipts).

\emph{Receipts}: $0$ violations in $27{,}000$ matrix elements; maximal one- and two-unit amplitudes $0.175384$ and $0.030759$, saturating $\sqrt{\varepsilon_{\rm eff}}$ and $\varepsilon_{\rm eff}$ to ten digits in the explicit model.

\textbf{Scope correction (withdrawn corollary).}  An earlier draft claimed "oscillations require the seam — the mechanism predicts its own mediator."  That conflated two gradings: if the Dirac sector is charged consistently, the claim's premise breaks the spectrum (Section 3); if the Dirac sector is outside the graded algebra, sectors connect without the seam and the corollary is vacuous.  Withdrawn.  What survives is the unglamorous statement: \emph{within the graded (sealed) sector}, charge steps cost seam insertions.

\section*{3. The obstruction by conjunction (central result)}
The natural realization — FN charges on the right-handed seals, $M_R$ hierarchical, Dirac sector neutral — \textbf{fails twice over}, and we record both failures as the paper's main content:

\textbf{(a) The textbook cancellation.}  If the Dirac couplings are charged consistently with the grading, the FN suppressions cancel identically in the seesaw, $m_\nu = v^2\,Y M_R^{-1}Y^{\mathsf T}$ — the well-known result that FN charges on singlet neutrinos drop out of the light masses (see e.g. [2]); and half-integer RH charges make the Dirac column amplitudes exactly zero under a unit-shifting seam, collapsing the rank.  An earlier draft evaded this by declaring the Dirac sector "neutral" while keeping the selection rule — an inconsistency the review exposed.

\textbf{(b) Parity, scoped, plus the scan.}  Majorana bilinears $\nu^c_i\nu^c_j$ carry total charge $x_i + x_j$.  For \textbf{integer} charges the diagonals carry even powers of $\lambda$ — a parity statement that by itself constrains only diagonal-dominant textures (off-diagonal-dominant blocks can carry odd-power eigenvalue magnitudes), so the integer-charge obstruction is completed by \emph{exhaustive enumeration}: over charge vectors in $\{0..3\}^3$ with unit $\pm$ coefficients on every rule-allowed entry, \textbf{zero} textures reach the rung spectrum.  The closest approach — metric pinned as the max over the two independent ratio deviations $|m_1/m_3 - \varepsilon_{\rm eff}|$ and $|m_2/m_3 - \sqrt{\varepsilon_{\rm eff}}|$, \textbf{computed on the seesaw-inverted (light) spectrum}: the charges sit on $M_R$ and the rung target is for the light masses, so $m_i \propto 1/\mu_i$ with $\mu_i$ the eigenvalue magnitudes of the scanned texture (equivalently, the two components are $|\mu_1/\mu_3 - \varepsilon_{\rm eff}|$ and $|\mu_1/\mu_2 - \sqrt{\varepsilon_{\rm eff}}|$) — is $0.014$, with the \textbf{witness published} so the figure is checkable from one matrix: $q = (3,2,3)$, coefficient pattern $\{11{:}{+}1, 12{:}{+}1, 13{:}{+}1, 22{:}{-}1, 23{:}{+}1, 33{:}0\}$ on magnitudes $\lambda^{q_i+q_j}$, component deviations $(0.0144, 0.0092)$. (Applying the metric to the heavy spectrum \emph{uninverted} gives $(0.0144, 0.0770)$ on this same witness; an earlier draft attributed that $0.077$ to "a less complete scan variant" — it is in fact the witness's own uninverted second component, and the inversion convention was unstated — the same class of convention bug a previous revision fixed for the charge box, caught in review and now pinned.)  The zero-exact-hits statement is robust under the pinned convention: a full independent rerun of the $\{0..3\}^3 \times \{0,\pm1\}^6$ scan confirms zero exact hits, best deviation $0.0144$.  Scope honesty: the charge box $\{0..3\}^3$ is a choice; entries beyond it carry magnitudes $\le \lambda^7 \approx 5\times 10^{-6}$ and push textures further from the target, but "exhaustive" means exhaustive-within-the-box.  For \textbf{half-odd} charges, parity does \emph{not} obstruct the Majorana side — $q = (1, \tfrac12, 0)$ makes $M_R = \mathrm{diag}(\lambda^2, \lambda, 1)$ rule-allowed and exactly on the rungs — but clause (a) kills the realization: the $q = \tfrac12$ Dirac column carries half-integer charge, balances no integer seam count, and is exactly zero, collapsing the light-mass rank.  \textbf{The obstruction is the conjunction of (a) and (b)-plus-scan}, and we state it as such; an earlier draft presented (b) as a standalone parity theorem (false for half the charge lattice) and left the scan's spectrum-inversion convention unstated (the source of the misattributed $0.077$ above) — both corrected here, with the scan scope and conventions now pinned.

\textbf{The surviving realization.}  Half-unit charges on the lepton \emph{doublets}, $q_L = (1, \tfrac12, 0)$, acting on the Weinberg operator $(L_iH)(L_jH)$ with entry charges $q_i + q_j$ and magnitudes $\lambda^{q_i+q_j}$ at integer insertions ($\lambda = \sqrt{\varepsilon_{\rm eff}}$): allowed entries are $(11), (22), (33), (13)$; $(12)$ and $(23)$ are parity-forbidden. This realization is \emph{selected, not chosen}, by two facts: (i) \textbf{integer} doublet charges on the Weinberg operator are excluded by the same exhaustive enumeration run on the \emph{direct} spectrum (no seesaw inversion here) — zero exact hits over the full $\{0..3\}^3 \times \{0,\pm1\}^6$ box, $30{,}624$ full-rank nondegenerate textures, best deviation $0.0770$, witness in the receipts; and (ii) $q_L = (1, \tfrac12, 0)$ is the \emph{unique} normalized assignment realizing the rung diagonal: $\mathrm{diag}(\lambda^{2q_1}, \lambda^{2q_2}, \lambda^{2q_3}) = \mathrm{diag}(\lambda^2, \lambda, 1)$ forces $2q_1 = 2$, $2q_2 = 1$, with the common shift fixed by $q_3 = 0$ (a shift rescales the spectrum, not its ratios). Convention, pinned: entries are $0$ or $\pm 1$ times the allowed magnitude; full rank and a nondegenerate ordering ($m_1 < m_2 < m_3$, all nonzero) required.  The exhaustive menu then has exactly \textbf{four} members (a previous revision said three, missing the quasi-degenerate fourth — caught in review; for a paper whose method is exhaustive enumeration, we print the complete list):

\begin{small}\begin{verbatim}
  texture            insertions  spectrum ratios       S
  diag(l^2, l, 1)        3       (e_eff, sqrt(e_eff), 1) 0.172747
  m13 + m22 + m33        2       (0.0290, 0.1703, 1)     0.167880
  all four, sgn(m11
   m33) = -1             4       (0.0581, 0.1704, 1)     0.160500
  m11 + m22 + m13,
   m33 = 0               4       (0.839, 0.916, 1)       0.675516
  (the third and fourth are already excluded by data - at 7.4
   and ~300 sigma respectively - and are listed as falsified
   members, not hidden; an earlier draft printed 0.169993 for
   the second member, a wrong closed-form shortcut - the correct
   value follows from the printed ratios themselves: S^2 =
   m1/(1+2m1), m1 = (sqrt(1+4*e_eff)-1)/2, receipted)
\end{verbatim}\end{small}
\section*{4. The two-point test}
Against JUNO's first data — $S = 0.17283 \pm 0.00167$, \textbf{our combination} of JUNO's $\Delta m^2_{21} = (7.50\pm 0.12)\times 10^{-5}$ [3] with the global-fit $\Delta m^2_{31} = (2.511 \pm 0.027)\times 10^{-3}$ [6], not a JUNO-only number:

\begin{small}\begin{verbatim}
  member               S          now        notes
  rung point           0.172747   -0.05 s    (reference)
  cheap texture        0.167880   -3.0  s    7.3 s from the rung
                                             point at JUNO-final
  third member         0.160500   -7.4  s    already excluded;
  fourth member        0.675516   excluded   listed, not hidden
  (undressed point     0.175472   +1.6  s    the registered claim of
   for comparison                            the companion note [VII])
\end{verbatim}\end{small}
The two data-live members are separated by $7.3\sigma$ at JUNO-final precision ($\sigma_S^{\rm final} \approx 0.00067$: the projected sub-percent splittings, $\sigma_{21} \approx 0.6\%$ and $\sigma_{31} \approx 0.5\%$ [5], propagated as half the quadrature sum of the relative errors times $S$): the menu is adjudicable this decade. Pre-assigned outcome bands at JUNO-final, all branches covered: data within $2\sigma$ of $0.172747$ — rung point survives, cheap texture dead; within $2\sigma$ of $0.167880$ — the reverse, and minimality scores a sharp win; \textbf{within $2\sigma_{\rm final}$ of $0.175472$ — both menu members die while [VII]'s undressed point survives}; anywhere else — the whole construction is dead and we say so.  Stated against interest: the program's own minimality principle (fewest insertions), applied as elsewhere in the corpus, favors the \textbf{cheap texture} — currently the \emph{worst-fitting live member at $-3.0\sigma$}.  If JUNO-final confirms today's central value, minimality-as-selector fails here and we will say so.

What the menu does \emph{not} include: mixing angles and CP phases. The allowed Weinberg textures are too sparse to generate the observed mixing from the neutrino sector alone (the $(12), (23)$ entries are parity-forbidden), so the PMNS matrix must come dominantly from the charged-lepton sector — the earlier draft's $\theta_{23}$ and $\delta_{CP}$ "predictions" are \textbf{withdrawn} (Section 7).  One genuine charged-sector \emph{constraint} the realization does impose, stated since review caught us denying it: $q_{L_2} = \tfrac12$ makes the muon Yukawa row $\bar L_2 e_R H$ carry half-integer charge, which vanishes unless the right-handed charged-lepton charges compensate modulo 1 — half-odd $e_{R_2}$ is an \textbf{additional named assumption} (a massless muon is not negotiable).  So the mechanism says nothing about charged \emph{masses} but does constrain charged \emph{charges}; the abstract's withdrawal item is scoped accordingly.

\section*{5. The one-line input, with its aggravating fact}
The mechanism's single quantitative input: the seam weight is the companion theory's first-order marginality, $p = \varepsilon = 3\kappa - 1$, with $\kappa$ in closed form ([C0]; no Standard-Model input anywhere in the chain).  A pre-registered referent scan (ten exact fixed-point candidates, 25 digits) found no independent quantity equal to $\varepsilon$: the identification is not derived, and we state the aggravating fact a referee located: $\varepsilon$ is the coefficient of a \emph{first-order} law whose sign prediction fails at the defining weight $w = 3$ (first-order $-0.0051$ vs exact $+0.0084$; [C0] itself warns against extrapolating its truncations).  The input is therefore: \emph{the seam weight is a specific coefficient of an expansion that is not trustworthy at the weight that motivates it.}  We keep the input because the resulting point sits at $-0.05\sigma$ and the falsification machinery is cheap; we decline to dress it up.

\section*{6. Chronology, trials, and look-elsewhere}
Disclosed in full.  (i) The \emph{value} $\varepsilon_{\rm eff} = \varepsilon(1-\varepsilon)$ was noticed \textbf{after} JUNO's first release; the program's registered claim is the \emph{undressed} point ($0.175472$ at exact $\varepsilon$; see [VII] and its precision correction), which that data disfavors at $1.6\sigma$.  (ii) The look-elsewhere cost was counted before this mechanism existed. The twelve forms, enumerated (base $\varepsilon\cdot f$): $f \in \{1-\varepsilon$, $1/(1+\varepsilon)$, $e^{-\varepsilon}$, $(1-\varepsilon)^2$, $1-2\varepsilon$, $1-\varepsilon/2$, $1/(1+2\varepsilon)$, $1-\varepsilon^2$, $\cos\varepsilon$, $1-\kappa\varepsilon$, $1-\theta\varepsilon$, $1-\eta\varepsilon\}$: six fit the current central value within $1\sigma$; three (including $\varepsilon_{\rm eff} = \varepsilon(1-\varepsilon)$) are mutually inseparable at $0.2\sigma$ even at JUNO-final — so the data can confirm or kill the one-factor \emph{class}, never select $\varepsilon(1-\varepsilon)$ within it, and a theorem pinning any cluster member would have looked equally prophetic.  The Seam Theorem's evidential weight is therefore \emph{zero until its premises are tested elsewhere}; its value is that it converts a free parameter into a falsifiable conjunction.  (iii) No claim in this paper is "registered" in any externally verifiable sense until the program's named deposit exists; the internal registration documents carry their own dating and say the same.

\section*{7. Withdrawn claims (kept visible)}
\begin{enumerate}
\item "$\sin^2\theta_{23} = 1/2$ exactly" — rested on requiring an \emph{orthogonal} neutral Dirac matrix, a measure-zero assumption with no symmetry behind it, which also smuggled CP conservation (real orthogonal) and silently fixed $\theta_{12}, \theta_{13}$ to excluded values.  Withdrawn as a mechanism prediction; what remains is a texture observation in the program's corpus (its registration document carries the same withdrawal).
\item "$J = 0$, $\delta_{CP} \in \{0, \pi\}$" — the phase-gauge argument covers the seam (Majorana) sector only; the Dirac and charged-lepton phases are untouched.  Withdrawn.
\item "Oscillations require the seam" — grading-scope conflation (Section 2).  Withdrawn.
\item The discrete $m_{\beta\beta}$ menu — computed with undisclosed measured mixing angles as inputs; meaningless once (1) is withdrawn.  Withdrawn.
\item "Exact rungs for any neutral Dirac sector" — exactness was an artifact of the orthogonality fiat; the honest exact statement is the Weinberg-realization menu of Section 3.
\item All charged-fermion \emph{mass} statements — the earlier draft quantified charged-sector "gaps" with an unreproducible metric; no such number survives.  (Not withdrawn: the charge-mod-1 compensation \emph{constraint} of Section 4, which the realization genuinely imposes and which carries its own named assumption.)
\end{enumerate}
\section*{8. Relation to Froggatt–Nielsen models}
The mechanism class is FN [1].  Prior art on \emph{derived} suppression exists and is engaged per our standing commitment: in anomalous- $U(1)$ (Green–Schwarz) FN models the parameter is computed from the Fayet–Iliopoulos term, $\lambda \approx 0.17$–$0.22$ (Binétruy–Ramond; Dudas–Pokorski–Savoy; Ibáñez–Ross [4]) — a range that, we note with appropriate discomfort, contains $\sqrt{\varepsilon_{\rm eff}} = 0.175$.  What is different here: the suppression is pinned by a stationarity theorem in a record formalism rather than by an anomaly coefficient, the charge's selection rule is exact rather than order-of-magnitude, and the realization question (this paper's obstruction and menu) is answered by exhaustive enumeration at unit coefficients.  What is \emph{not} different: like every FN model, the O(1)-coefficient freedom is the mechanism's soft underbelly, and our unit- coefficient discipline is a named choice, not a derivation.

\section*{Reproducibility}
Fixed-seed scripts regenerate every number (the seam lemmas, the selection-rule receipts, the parity-obstruction scan, the Weinberg menu, the prediction table); bit-identical reruns; code and canonical outputs accompany the submission.

\section*{References}
[1] C. D. Froggatt, H. B. Nielsen, \emph{Nucl. Phys. B} \textbf{147} (1979) 277.

[2] G. Altarelli, F. Feruglio, Models of neutrino masses and mixings, \emph{New J. Phys.} \textbf{6} (2004) 106 — including the cancellation of singlet-neutrino FN charges in the seesaw.

[3] JUNO collaboration, arXiv:2511.14593; \emph{Nature} \textbf{654} (2026) 343 — first measurement, $\Delta m^2_{21} = (7.50 \pm 0.12)\times 10^{-5}\,\mathrm{eV}^2$.

[4] P. Binétruy, P. Ramond, \emph{Phys. Lett. B} \textbf{350} (1995) 49; E. Dudas, S. Pokorski, C. A. Savoy, \emph{Phys. Lett. B} \textbf{356} (1995) 45; L. E. Ibáñez, G. G. Ross, \emph{Phys. Lett. B} \textbf{332} (1994) 100 — anomalous-$U(1)$ FN with FI-derived $\lambda$.

[5] JUNO collaboration, JUNO physics and detector, \emph{Prog. Part. Nucl. Phys.} \textbf{123} (2022) 103927 — projected sub-percent splitting precisions (an earlier draft cited the DUNE TDR and Hyper-Kamiokande design report here; those experiments do not measure $\Delta m^2_{21}$ at sub-percent — misattribution, corrected).

[6] I. Esteban et al., NuFIT — global-fit $\Delta m^2_{31}$ input; version sensitivity applies to the hybrid $S$ above.

\textbf{Companions:} [C0] the fixed-point theory (the constant's pedigree); [C2] the representation theorem; [VII] the registered undressed point and its correction record (the two claims are mutually exclusive and cross-declared).

\end{document}
